% !PRINTED 645
\weberpage{645}

But if $m$ is a power of 2, the congruence
\[
p^h \equiv -1 \ (\mathrm{mod}\ m)
\]
is possible only when $p \equiv -1\ (\mathrm{mod}\ m)$. For every even power of an odd number is $\equiv + 1\ (\mathrm{mod}\ 8)$ and therefore cannot be $\equiv -1\ (\mathrm{mod}\ m)$. But if $h$ is odd, then
\[
\frac{p^h + 1}{p + 1} = p^{h-1} - p^{h-2} + \cdots + 1
\]
is itself odd, and consequently $p^h + 1$ is divisible by no higher power of 2 than $p + 1$. Accordingly we may add to Theorem 2. the following closer determination:

3. \gesperrt{If $m$ is odd, the primes $p_1$ belong to an even exponent, the primes $p_2$ to an odd exponent.}

\gesperrt{If $m$ is even, the primes of the form $km - 1$ belong to the first kind, and all other odd primes to the second.}

\begin{center}
\S.\,178.\\[4pt]
\gesperrt{The Units of the Field $H_m$.}
\end{center}

The units belonging to the field $H_m$ are of particular importance for the applications. The units of the field $\Omega_m$ can also be reduced to them according to the following theorem:

1. \gesperrt{Every unit of the field $\Omega_m$ is the product of a root of unity and a real unit of the field $\Omega_m$.}

If we denote any unit of the field $\Omega_m$ by $\mathfrak{E}(r)$, then $\mathfrak{E}(r^{-1})$ is the conjugate imaginary value, which likewise represents a unit, and the quotient $\mathfrak{E}(r) : \mathfrak{E}(r^{-1})$ is again a unit, whose absolute value
\weq{1}{\sqrt{\dfrac{\mathfrak{E}(r)}{\mathfrak{E}(r^{-1})}\ \dfrac{\mathfrak{E}(r^{-1})}{\mathfrak{E}(r)}}}
is equal to 1, and consequently this quotient is a root of unity and hence $= \pm r^h$, where $h$ is any integer exponent (§. 175, Theorem 1., 2.). We therefore have
\weq{2}{\mathfrak{E}(r) = \pm r^h \mathfrak{E}(r^{-1}).}
